% ============================================================================
% Fichier  : partie4/P04_C22_tests_intrusion_code.tex
% Titre    : Tests d'Intrusion et Audit de Code
% Auteur   : MINKA MI NGUIDJOÏ Thierry Emmanuel
% Version  : 1.0 -- novembre 2026
% ============================================================================
% Chapitre propre à INF3226 (ACSI C16 ne traite que la note d'autorisation).
% Sources : loi n° 2010/012, art. 65 (2), 69, 70 et 86 (1), texte officiel
% lu le 06/11/2026 (fichier loi_2010-012_cybersecurite_cybercriminalite.pdf) ;
% NIST SP 800-115 (phases d'un test) ; OWASP Top 10 (2021) ; OWASP ASVS
% 4.0.3 (2021) ; OWASP Web Security Testing Guide ; FIRST CVSS v3.1.
% Faits nouveaux (LiZbiZ) : portail de suivi des expéditions exposé en 443 ;
% test d'intrusion en boîte grise du 8 au 12 juin 2026 ; 1 critique
% (injection SQL), 2 élevées, 4 moyennes, 3 faibles. Base du lab
% ASI-LAB-12 (cible vulnérable isolée).
% Corrigés : corriges/Corr_C22.tex.
% ============================================================================

\chapter{Tests d'Intrusion et Audit de Code}
\label{chap:intrusion}
\chaptermark{Intrusion et audit de code}

\epigraph{%
    « Le testeur et l'attaquant font les mêmes gestes.\\
    Seule l'autorisation écrite les sépare. »%
}{Règle des tests d'intrusion}

\niveauChapitre{Expert}{%
    Maîtrise du Chapitre~\ref{chap:audit-technique} (infrastructure) et du
    Chapitre~\ref{chap:droit} (loi n°~2010/012)~; notions de
    programmation web%
}

\begin{boxobjectifs}
\begin{itemize}
    \item Distinguer le scan de vulnérabilités, le test d'intrusion et
          l'audit de code.
    \item Connaître les articles de la loi n°~2010/012 qui rendent
          l'autorisation écrite indispensable, et rédiger une note
          d'autorisation.
    \item Décrire les phases d'un test d'intrusion et ses modalités (boîte
          noire, grise, blanche).
    \item Utiliser l'OWASP Top~10 et l'OWASP ASVS comme référentiels
          d'audit d'une application web.
    \item Distinguer l'analyse statique, l'analyse dynamique et la revue
          manuelle du code.
    \item Lire et auditer un rapport de test d'intrusion, et en tirer les
          cotes des fiches d'assertion.
\end{itemize}
\end{boxobjectifs}

\bigskip

% ============================================================================
\section*{Question directrice : un attaquant pourrait-il entrer~?}
% ============================================================================

LiZbiZ expose sur Internet un portail qui permet à ses clients de suivre
leurs expéditions. L'audit des configurations (Chapitre~\ref{chap:audit-technique})
a montré des règles de pare-feu permissives dans deux centres. Le conseil
pose une question plus directe~: un attaquant extérieur pourrait-il, par ce
portail, atteindre les données des clients~?

Pour répondre, il faut tenter de le faire, comme un attaquant, mais avec
une autorisation, un périmètre et des règles. C'est l'objet du test
d'intrusion.

\medskip
\textbf{Plan du chapitre.}
\begin{enumerate}
    \item Trois techniques (section~\ref{sec:c22-trois}).
    \item Le cadre légal et l'autorisation (section~\ref{sec:c22-droit}).
    \item Conduire un test d'intrusion (section~\ref{sec:c22-conduire}).
    \item Les référentiels de l'OWASP (section~\ref{sec:c22-owasp}).
    \item L'audit de code (section~\ref{sec:c22-code}).
    \item Le rapport et son audit (section~\ref{sec:c22-rapport}).
\end{enumerate}

\separateur

% ============================================================================
\section{Trois Techniques}
\label{sec:c22-trois}
% ============================================================================

\begin{tableaucours}{L{2.8cm}L{4.4cm}L{2.6cm}L{2.4cm}}{%
    \tableauHeader{Technique} &
    \tableauHeader{Principe} &
    \tableauHeader{Ce qu'elle établit} &
    \tableauHeader{Limite}
}
Scan de vulnérabilités & Outil automatique qui compare les versions et
configurations à une base de failles connues & Une liste de faiblesses
probables & Faux positifs~; ne prouve pas l'exploitation \\
Test d'intrusion & Un testeur exploite les faiblesses, les enchaîne, et
mesure jusqu'où il va & Un chemin d'attaque démontré & Ponctuel~; limité
au périmètre et au temps \\
Audit de code & Lecture du code source, outillée ou manuelle & Les
défauts à la source, y compris ceux qu'un test ne voit pas & Demande le
code et des compétences de développeur \\
\end{tableaucours}
\vspace{-0.8em}
\begin{center}{\small\itshape Tableau~22.1. Scan, test d'intrusion, audit de
code}\end{center}

Les trois se complètent~: le scan trie, le test démontre, l'audit de code
explique et corrige à la source. Le guide du NIST sur l'évaluation
technique de la sécurité les classe parmi les techniques de revue,
d'identification des cibles et de validation des
vulnérabilités~\cite{nist2008sp800115}.

% ============================================================================
\section{Le Cadre Légal et l'Autorisation}
\label{sec:c22-droit}
% ============================================================================

\subsection{Ce que la loi réprime}

Un testeur accomplit matériellement des actes que la loi réprime
lorsqu'ils sont faits sans droit.

\begin{boxjuris}
\textbf{\loicam{n°~2010/012}} (texte officiel)

\medskip
\textbf{Article~65, alinéa~2}~: \textit{« Est puni des peines prévues à
l'alinéa 1 ci-dessus, tout accès non autorisé, à l'ensemble ou à une
partie d'un réseau de communications électroniques ou d'un système
d'information ou d'un équipement terminal. »}

\medskip
\textbf{Article~69}~: \textit{« Est puni d'un emprisonnement de cinq (05) à
dix (10) ans et d'une amende de 10.000.000 (dix millions) à 100.000.000
(cent millions) F CFA ou de l'une de ces deux peines seulement, celui qui
accède sans droit, et en violation des mesures de sécurité, à l'ensemble
ou à une partie d'un réseau de communications électroniques, d'un système
d'information ou d'un équipement terminal, afin d'obtenir des informations
ou des données, en relation avec un système d'information connecté à un
autre système d'information. »}

\medskip
\textbf{Article~70}~: \textit{« Est puni d'une amende de 1.000.000 (un
million) à 5.000.000 (cinq millions) F CFA, celui qui provoque par
saturation, l'attaque d'une ressource de réseau de communications
électroniques ou d'un système d'information dans le but de l'effondrer en
empêchant la réalisation des services attendus. »}

\medskip
\textbf{Article~86, alinéa~1}~: \textit{« Est puni des peines prévues
l'article 71 ci-dessus, celui qui importe, détient, offre, cède, vend ou
met à disposition, sous quelle que forme que ce soit, un programme
informatique, un mot de passe, un code d'accès ou toutes données
informatiques similaires conçus et ou spécialement adaptés, pour permettre
d'accéder, à tout ou partie d'un réseau de communications électroniques ou
d'un système d'information. »}

\medskip
\textbf{Analyse.} L'autorisation écrite du responsable du système est ce
qui donne au testeur le « droit » d'accéder~: sans elle, l'accès tombe sous
l'article~65 ou~69. Un test de charge non encadré peut produire l'effet visé
par l'article~70. L'article~86 vise la détention d'outils « spécialement
adaptés » pour accéder~: le cabinet documente l'usage professionnel de ses
outils (mission, autorisation) et les conserve sous contrôle. La portée
exacte de ces articles pour un professionnel autorisé relève de
l'interprétation du juge~; la prudence impose de ne jamais agir sans
autorisation écrite.
\end{boxjuris}

\subsection{La note d'autorisation}

\begin{boxprocedure}
\textbf{Contenu de la note d'autorisation d'un test technique}
\begin{enumerate}
    \item \textbf{Signataire}~: une personne habilitée à engager
          l'organisation sur les systèmes visés, et, pour un système
          hébergé, l'accord écrit de l'hébergeur.
    \item \textbf{Périmètre}~: adresses, systèmes et applications inclus~;
          systèmes explicitement exclus.
    \item \textbf{Méthodes}~: lecture de configuration, scan de ports, scan
          de vulnérabilités, test d'intrusion~; ce qui est interdit (déni
          de service, modification ou suppression de données, accès aux
          données réelles au-delà de la preuve).
    \item \textbf{Fenêtres}~: dates et heures autorisées, hors des périodes
          critiques.
    \item \textbf{Contacts}~: un interlocuteur joignable pendant les tests,
          et une procédure d'arrêt immédiat.
    \item \textbf{Résultats}~: confidentialité, conservation, destruction
          à la fin de la mission (secret des experts, art.~14).
\end{enumerate}
\end{boxprocedure}

\begin{boxalert}
\textbf{Violation typique.} Lancer un scan « rapide » du réseau de Garoua
depuis le portable de l'auditeur, la veille de la signature de la note
d'autorisation, « pour gagner du temps ».

\smallskip
\textbf{Règle~:} aucun test avant la signature, aucun test hors du
périmètre ni de la fenêtre. Un constat obtenu sans droit n'est pas
opposable, et il expose l'auditeur.
\end{boxalert}

% ============================================================================
\section{Conduire un Test d'Intrusion}
\label{sec:c22-conduire}
% ============================================================================

\subsection{Quatre phases}

Le guide du NIST décrit un test en quatre phases~\cite{nist2008sp800115}.

\begin{tableaucours}{L{2.8cm}L{5.0cm}L{4.4cm}}{%
    \tableauHeader{Phase} &
    \tableauHeader{Activités} &
    \tableauHeader{Trace à conserver}
}
Planification & Périmètre, règles d'engagement, autorisation, contacts &
Note d'autorisation signée \\
Découverte & Recensement des cibles, des services, des versions~; analyse
des vulnérabilités & Sorties d'outils brutes, horodatées, avec empreintes \\
Attaque & Exploitation contrôlée, élévation de privilèges, rebond vers
d'autres systèmes, dans le périmètre & Journal de bord minute par minute,
captures \\
Rapport & Constats, preuves, gravité, recommandations & Rapport, annexes,
destruction des données obtenues \\
\end{tableaucours}
\vspace{-0.8em}
\begin{center}{\small\itshape Tableau~22.2. Les phases d'un test d'intrusion
(d'après le NIST SP~800-115)}\end{center}

\subsection{Boîte noire, grise ou blanche}

\begin{tableaucours}{L{2.4cm}L{5.2cm}L{4.6cm}}{%
    \tableauHeader{Modalité} &
    \tableauHeader{Ce que sait le testeur} &
    \tableauHeader{Usage}
}
Boîte noire & Rien, sinon l'adresse publique & Simuler un attaquant
extérieur sans information \\
Boîte grise & Un compte utilisateur, la documentation fonctionnelle &
Simuler un client ou un agent malveillant~; le meilleur rapport
coût-couverture \\
Boîte blanche & Le code, l'architecture, les comptes & Couverture maximale,
proche de l'audit de code \\
\end{tableaucours}
\vspace{-0.8em}
\begin{center}{\small\itshape Tableau~22.3. Les modalités d'un test}\end{center}

\begin{boxscenario}{Opération LIZBIZ. Le test du portail de suivi}
\textbf{Cadre.} Test en boîte grise du portail de suivi des expéditions
(un compte client de test), du 8 au 12~juin~2026, de 20~h à 6~h, sur une
copie de production isolée~; note signée par le directeur général et par
l'hébergeur~; déni de service exclu.

\medskip
\textbf{Constatations} (rapport du testeur, preuves en annexe).
\begin{itemize}
    \item \textbf{Critique}~: injection SQL dans le paramètre de numéro de
          suivi~; extraction démontrée de dix enregistrements fictifs de
          la table des clients (CVSS~9,8).
    \item \textbf{Élevées} (2)~: un client peut afficher l'expédition d'un
          autre en modifiant l'identifiant dans l'adresse~; mots de passe
          stockés avec une fonction de hachage obsolète, sans sel.
    \item \textbf{Moyennes} (4)~: en-têtes de sécurité absents, messages
          d'erreur détaillés, cookie de session sans attribut sécurisé,
          bibliothèque JavaScript obsolète.
    \item \textbf{Faibles} (3)~: informations de version exposées.
\end{itemize}

\textbf{Analyse} (confiance élevée). La réponse au conseil est oui~: un
attaquant extérieur avec un compte client gratuit peut atteindre les
données des clients. Le flux de la DMZ vers la base, ouvert à Douala
(Chapitre~\ref{chap:audit-technique}), aurait permis d'aller plus loin~;
le testeur s'est arrêté à la preuve, conformément à la note.
\end{boxscenario}

% ============================================================================
\section{Les Référentiels de l'OWASP}
\label{sec:c22-owasp}
% ============================================================================

\subsection{L'OWASP Top 10}

L'OWASP (\textit{Open Worldwide Application Security Project}) publie
une liste des dix catégories de risques les plus critiques pour les
applications web~\cite{owasp2021top10}.

\begin{tableaucours}{cL{5.0cm}L{6.0cm}}{%
    \tableauHeader{Rang} &
    \tableauHeader{Catégorie (2021)} &
    \tableauHeader{Exemple au portail de LiZbiZ}
}
A01 & Contrôle d'accès défaillant & Voir l'expédition d'un autre client \\
A02 & Défaillances cryptographiques & Hachage obsolète sans sel \\
A03 & Injection & Injection SQL sur le numéro de suivi \\
A04 & Conception non sécurisée & Absence de modèle de menace \\
A05 & Mauvaise configuration de sécurité & Messages d'erreur détaillés \\
A06 & Composants vulnérables et obsolètes & Bibliothèque JavaScript
obsolète \\
A07 & Identification et authentification défaillantes & Pas de
verrouillage après échecs \\
A08 & Intégrité des logiciels et des données & Mises à jour non signées \\
A09 & Journalisation et surveillance défaillantes & Tentatives d'injection
non détectées \\
A10 & Falsification de requête côté serveur & (non constaté) \\
\end{tableaucours}
\vspace{-0.8em}
\begin{center}{\small\itshape Tableau~22.4. L'OWASP Top~10 (2021) et le
portail de LiZbiZ}\end{center}

\begin{boxinfo}
\textbf{Le Top~10 n'est pas une norme.} C'est un document de
sensibilisation. Pour un audit, on utilise un référentiel d'exigences
vérifiables~: l'ASVS.
\end{boxinfo}

\subsection{L'OWASP ASVS}

L'\emph{Application Security Verification Standard} (version 4.0.3)
énumère plusieurs centaines d'exigences de sécurité vérifiables, réparties
en chapitres (architecture, authentification, sessions, contrôle d'accès,
validation des entrées, cryptographie, journalisation\ldots) et en trois
niveaux~\cite{owasp2021asvs}.

\begin{tableaucours}{cL{4.6cm}L{6.4cm}}{%
    \tableauHeader{Niveau} &
    \tableauHeader{Pour qui} &
    \tableauHeader{Vérification}
}
1 & Toute application & Exigences vérifiables par des tests extérieurs \\
2 & Applications qui traitent des données sensibles & La plupart des
exigences, dont celles qui demandent l'accès au code ou à la conception \\
3 & Applications critiques (paiement, santé, infrastructures) & Toutes les
exigences, revue approfondie \\
\end{tableaucours}
\vspace{-0.8em}
\begin{center}{\small\itshape Tableau~22.5. Les niveaux de l'OWASP
ASVS}\end{center}

Pour l'auditeur, l'ASVS fournit des critères atomiques, prêts à entrer dans
une fiche d'assertion~: « toutes les requêtes à la base utilisent des
requêtes paramétrées » est une exigence de l'ASVS, et un critère.

% ============================================================================
\section{L'Audit de Code}
\label{sec:c22-code}
% ============================================================================

\subsection{Trois approches}

\begin{tableaucours}{L{3.0cm}L{4.8cm}L{4.4cm}}{%
    \tableauHeader{Approche} &
    \tableauHeader{Principe} &
    \tableauHeader{Forces et limites}
}
Analyse statique (SAST) & Un outil lit le code sans l'exécuter et signale
les motifs dangereux & Exhaustive, rapide~; beaucoup de faux positifs \\
Analyse dynamique (DAST) & Un outil teste l'application en fonctionnement,
de l'extérieur & Trouve ce qui est exploitable~; ne dit pas où est le
défaut \\
Revue manuelle & Un développeur expérimenté lit les parties sensibles
(authentification, accès, requêtes) & Comprend la logique métier~; coûteuse \\
\end{tableaucours}
\vspace{-0.8em}
\begin{center}{\small\itshape Tableau~22.6. Trois approches de l'audit de
code}\end{center}

\subsection{Ce que montre le code}

\begin{boxoutils}{L'injection du portail, à la source}
\begin{lstlisting}[language=Python]
# code vulnérable : la saisie est concaténée à la requête
def suivi(numero):
    sql = "SELECT * FROM expeditions WHERE numero = '" + numero + "'"
    return db.execute(sql)

# code corrigé : requête paramétrée, la saisie n'est jamais du SQL
def suivi(numero, client_id):
    sql = "SELECT * FROM expeditions WHERE numero = ? AND client_id = ?"
    return db.execute(sql, (numero, client_id))
\end{lstlisting}
La correction traite deux constats à la fois~: l'injection (A03), par la
requête paramétrée, et le contrôle d'accès (A01), par le filtre sur le
client connecté.
\end{boxoutils}

\begin{boxalert}
\textbf{Violation typique.} Corriger l'injection en filtrant les apostrophes
dans le numéro de suivi.

\smallskip
\textbf{Règle~:} une liste de caractères interdits se contourne. La
correction attendue est structurelle (requête paramétrée), et l'auditeur la
vérifie dans le code, pas seulement par un nouveau test.
\end{boxalert}

% ============================================================================
\section{Le Rapport et son Audit}
\label{sec:c22-rapport}
% ============================================================================

\subsection{Ce qu'un rapport doit contenir}

\begin{boxprocedure}
\textbf{Le rapport de test d'intrusion}
\begin{enumerate}
    \item Une synthèse pour la direction~: ce qu'un attaquant peut faire,
          en une page.
    \item Le cadre~: autorisation, périmètre, modalité, fenêtres, outils et
          versions.
    \item Les constats, classés par gravité (CVSS), chacun avec sa preuve
          (requête, réponse, capture), ses conditions de reproduction et sa
          recommandation.
    \item Les chemins d'attaque~: comment les faiblesses s'enchaînent.
    \item Ce qui n'a pas été testé, et pourquoi.
    \item L'attestation de destruction des données obtenues.
\end{enumerate}
\end{boxprocedure}

\subsection{Du rapport aux fiches d'assertion}

L'auditeur SI qui reçoit un rapport de test d'intrusion ne le recopie pas.
Il en tire des cotes pour les fiches d'assertion concernées. Une faille
critique démontrée est un \NON{} établi sur le contrôle correspondant~;
l'absence de faille trouvée n'est pas un \OUI{}~: c'est l'absence de preuve
contraire dans un périmètre et un temps limités.

\begin{boxlecture}{Le portail de LiZbiZ après le test}
AUDIX CONSEIL cote huit exigences de l'ASVS (niveau~2) sur le portail, à
partir du rapport et d'une revue du code des fonctions de suivi.

\medskip
\textbf{Constatation.} 3~\OUI{} (en-têtes de transport chiffré, séparation
des environnements, absence de compte par défaut), 4~\NON{} (requêtes
paramétrées, contrôle d'accès par objet, stockage des mots de passe,
composants à jour), 1~\NEK{} (journalisation des attaques, aucun journal
remis)~: $\mu \approx 0{,}38$~; $U \approx 0{,}13$~; $H \approx 1{,}41$~bit.

\medskip
\textbf{Analyse} (confiance élevée). Non-conformité claire~: le portail ne
satisfait pas les exigences de base. La priorité est la faille critique,
corrigée à la source, puis vérifiée par un nouveau test.

\medskip
\hfill\lecture{0{,}38}{0{,}13}{1{,}41}
\end{boxlecture}

\begin{boxalert}
\textbf{Violation typique.} Conclure « l'application est sécurisée » parce
que le test n'a rien trouvé.

\smallskip
\textbf{Règle~:} un test ne trouve que ce qu'il a cherché, dans le temps
qu'il avait. Le rapport dit « aucune vulnérabilité exploitable n'a été
trouvée dans le périmètre et la durée du test », jamais « l'application est
sûre ».
\end{boxalert}

\begin{boxmetier}{Auditeur de sécurité}
\textbf{Conduire un test d'intrusion}
\begin{enumerate}
    \item Ne rien faire avant la note signée~; la garder sur soi pendant
          le test.
    \item Tenir un journal de bord minute par minute.
    \item S'arrêter à la preuve~: démontrer l'accès, ne pas exploiter les
          données.
    \item Signaler immédiatement une faille critique au contact désigné,
          sans attendre le rapport.
    \item Détruire les données obtenues et l'attester.
\end{enumerate}
\end{boxmetier}

\begin{boxmetier}{RSSI}
\textbf{Faire tester sans subir.} Le RSSI prévient la surveillance
(pour qu'un test autorisé ne déclenche pas une réponse à incident), prévoit
une copie de production lorsque c'est possible, fait corriger à la source
et commande un test de vérification après correction.
\end{boxmetier}

\separateur

\begin{boxresume}
\textbf{Ce qu'il faut retenir de ce chapitre~:}
\begin{itemize}
    \item \textbf{Scan, test d'intrusion, audit de code}~: trier, démontrer,
          expliquer à la source.
    \item \textbf{La loi n°~2010/012} (art.~65, 69, 70, 86) rend
          l'autorisation écrite indispensable~; aucun test sans elle,
          aucun test hors périmètre.
    \item \textbf{Un test} suit quatre phases (planification, découverte,
          attaque, rapport), en boîte noire, grise ou blanche.
    \item \textbf{L'OWASP Top~10} sensibilise~; \textbf{l'ASVS} fournit
          des exigences vérifiables, en trois niveaux.
    \item \textbf{L'audit de code} combine analyse statique, dynamique et
          revue manuelle~; une correction se vérifie à la source.
    \item \textbf{Un rapport de test} donne des \NON{} établis~; l'absence
          de faille trouvée n'est pas un \OUI{}.
\end{itemize}
\end{boxresume}

\begin{boxinfo}
\textbf{Points de contrôle de l'élève.} Je sais~: citer les articles 65,
69, 70 et 86 et leur rôle~; rédiger une note d'autorisation~; nommer les
quatre phases et les trois modalités~; associer un constat à une catégorie
du Top~10~; choisir un niveau ASVS~; corriger une injection par une
requête paramétrée~; tirer des cotes d'un rapport.
\end{boxinfo}

\bigskip

\begin{boxexo}[title={\ding{45}\quad Exercice~22.1. Quelle technique~? \niveaubase}]
Pour chaque besoin, choisissez la technique (Tableau~22.1)~: (a) savoir si
les 300~serveurs de la banque ont des versions vulnérables~; (b) savoir si
un employé malveillant peut devenir administrateur du domaine~; (c) savoir
pourquoi le module de paiement accepte des montants négatifs.
\end{boxexo}

\begin{boxexo}[title={\ding{45}\quad Exercice~22.2. La note d'autorisation \niveaumoyen}]
AUDIX CONSEIL doit réaliser un scan de vulnérabilités des quatre centres de
LiZbiZ. Le centre de Garoua est hébergé par un prestataire local.
\begin{enumerate}[(a)]
    \item Rédigez la note d'autorisation.
    \item Qui doit la signer pour Garoua, et pourquoi~?
    \item Le DSI propose d'ajouter un « test de charge » de l'application.
          Que répondez-vous, au regard de l'article~70~?
\end{enumerate}
\end{boxexo}

\begin{boxexo}[title={\ding{45}\quad Exercice~22.3. Boîte noire ou grise~? \niveaubase}]
Le conseil de la BANQUE SANAGA hésite entre un test en boîte noire de
l'application mobile et un test en boîte grise. Avec un budget de dix
jours, que recommandez-vous, et pourquoi~?
\end{boxexo}

\begin{boxexo}[title={\ding{45}\quad Exercice~22.4. Classer selon le Top 10 \niveaubase}]
Classez selon l'OWASP Top~10~: (a) une page d'administration accessible sans
connexion~; (b) une version de serveur web vieille de cinq ans~; (c)
l'absence de journal des connexions échouées~; (d) un mot de passe
transmis en HTTP.
\end{boxexo}

\begin{boxexo}[title={\ding{45}\quad Exercice~22.5. Une exigence ASVS \niveaumoyen}]
Transformez l'exigence « les requêtes à la base utilisent des requêtes
paramétrées » en fiche d'assertion complète pour le portail de LiZbiZ~:
objectif, critère, question, risque, conséquence, test, preuve, trace.
\end{boxexo}

\begin{boxexo}[title={\ding{45}\quad Exercice~22.6. Corriger le code \niveaumoyen}]
Une fonction de connexion construit la requête
\texttt{"SELECT * FROM users WHERE login='" + l + "' AND pwd='" + p + "'"}.
Expliquez l'attaque possible, et écrivez la version corrigée (requête
paramétrée, et vérification d'un hachage de mot de passe plutôt qu'une
comparaison en clair).
\end{boxexo}

\begin{boxexo}[title={\ding{45}\quad Exercice~22.7. Lire un rapport \niveaumoyen}]
Un rapport de test conclut~: « Aucune vulnérabilité critique. L'application
est sécurisée. » Le rapport ne mentionne ni le périmètre, ni la durée, ni
les comptes utilisés. Que vaut-il pour l'auditeur SI~? Quelles cotes en
tirez-vous~?
\end{boxexo}

\begin{boxexo}[title={\ding{45}\quad Exercice~22.8. Recalculer la lecture \niveaumoyen}]
Recalculez $\mu$, $U$ et $H$ de la lecture du portail (3/4/1). Après
correction, un nouveau test et la remise des journaux, la lecture devient
6/1/1. Recalculez, situez dans le quadrant, et dites ce qui reste à faire.
\end{boxexo}

\begin{boxexo}[title={\ding{45}\quad Exercice~22.9. La faille trouvée la nuit \niveauexpert}]
À 2~h du matin, pendant le test, le testeur découvre que l'injection SQL
permet aussi de lire la table des paiements \emph{de production}, alors
que la note limitait le test à une copie isolée~: la copie était mal
isolée. Que doit-il faire, dans quel ordre, et que contient sa
communication au contact désigné~? Quelles sont les conséquences pour la
suite de la mission~?
\end{boxexo}

\begin{boxexo}[title={\ding{45}\quad Exercice~22.10. Un programme de tests \niveauexpert}]
Proposez à la BANQUE SANAGA un programme annuel de tests techniques
(scans, tests d'intrusion, audit de code) pour son application mobile, son
site public et son système central, avec la fréquence, la modalité, le
niveau ASVS visé et la manière d'utiliser les résultats dans l'audit
interne.
\end{boxexo}

\bigskip

\begin{boxinfo}
\textbf{Pour aller plus loin.}
\begin{itemize}
    \item Loi n°~2010/012~\cite{loi2010012}, articles~65 à~72 et~86.
    \item NIST, SP~800-115~\cite{nist2008sp800115}.
    \item OWASP, \textit{Top 10:2021}~\cite{owasp2021top10} et
          \textit{ASVS 4.0.3}~\cite{owasp2021asvs}.
\end{itemize}
\end{boxinfo}

\begin{boxinfo}
\textbf{Transition.} Une infrastructure sûre doit aussi durer. Le
Chapitre~\ref{chap:audit-continuite} audite la production et la
continuité~: exploitation, sauvegardes, restaurations, plans de
continuité.
\end{boxinfo}
